\subsection{STM32F4xx Characteristics}
\label{chap:Back_ARM_STM32F4xx}

The present thesis is centered in the emulation of the STM32F405's I2C protocol. 
In this section, will be listed some of the STM32F405 features, extracted directly 
from its datasheet\cite{STM32F405Datasheet2026}\cite{STM32CortexM42026}. This microcontroller is a high-performance 32-bit 
microcontroller from STMicroelectronics, built around an Arm Cortex-M4 core with FPU, 
and DSP extensions, and an adaptive real-time accelerator that enables zero-wait-state 
code execution from Flash at clock frequencies up to 168 MHz. It delivers up 
to around 210 DMIPS and includes a memory protection unit, making it suitable for 
complex real-time embedded applications that require both high performance and 
robust software isolation.

In terms of memory, the STM32F405 integrates up to 1 MByte of embedded Flash 
and up to 196 KBytes of SRAM, including a 64 KByte core-coupled memory (CCM) 
region optimized for fast data access, along with OTP memory and a flexible static 
memory controller for external SRAM, PSRAM, NOR, NAND, and CompactFlash 
devices. The device operates from 1.8 V to 3.6 V, offers multiple clock sources 
(external crystal, internal 16 MHz RC, and 32 kHz oscillators for RTC), and supports 
several low-power modes (Sleep, Stop, Standby) with a dedicated VBAT domain for 
RTC and backup registers, enabling energy-efficient designs that retain timekeeping 
and critical data across power cycles.

The STM32F405 features rich analog capabilities, including three 12-bit ADCs 
with multi-channel and interleaved modes reaching multi-MSPS aggregate sampling 
rates, as well as two 12-bit DACs for signal generation. It also provides a 
16-stream general-purpose DMA controller with FIFO and burst support, a CRC unit, 
a true random number generator, and a large timer subsystem with up to 17 timers 
(16- and 32-bit) supporting capture/compare, PWM generation, pulse counting, 
and quadrature encoder interfaces, making it well suited for motor control, 
power conversion, and advanced control applications.

On the digital I/O and connectivity side, the STM32F405 exposes up to 140 
GPIOs (most of them fast and 5 V tolerant) and supports a wide range of communication 
peripherals: up to three I2C controllers, multiple USART/UART interfaces 
with support for protocols like LIN, IrDA, and ISO 7816, up to three SPI/I2S blocks, 
two CAN 2.0B controllers, SDIO, and a parallel camera interface. Advanced connectivity 
is further enhanced by an integrated Ethernet MAC with IEEE 1588v2 
hardware timestamping and DMA, as well as full-speed and high-speed USB 2.0 
OTG controllers with on-chip PHY and dedicated DMA, enabling the microcontroller 
to serve as a communication hub in networked and multimedia embedded 
systems.

For development and system integration, the MCU includes a serial-wire/JTAG 
debug interface, an Embedded Trace Macrocell for instruction tracing, an LCD 
parallel interface, a real-time clock with calendar and subseconds, and a 96-bit 
unique device ID that can be used for traceability or security-related features.